\chapter{Introduction}
\label{sec:introduction}

\section{Problem and Motivation}

Modern fleet monitoring depends on telemetry streams that describe vehicle
state, motion, energy availability, and accumulated usage. In electric
vehicles, measurements such as remaining energy, remaining range proxy, speed,
acceleration, mileage, and travelled distance are not independent facts. Their
operational meaning depends on how they relate: acceleration changes speed,
speed contributes to distance, driving behaviour affects energy consumption,
and remaining energy constrains how a vehicle can be interpreted in a fleet
context.

This relational structure makes graph-based telemetry analysis attractive. A
directed graph can encode proposed dependencies between telemetry features,
while an adjacency matrix makes those dependencies measurable and comparable.
Such representations can support inspection, diagnostics, knowledge-graph
construction, and downstream decision support. They are also difficult to
obtain automatically. Classical causal discovery requires assumptions about
data-generating processes, sufficient observations, and domain validation.
Manual graph construction by experts is more interpretable, but it is costly
and does not scale easily across vehicle systems, datasets, and use cases.

Large language models (LLMs) offer a different kind of support. Rather than
estimating a causal model directly from raw data, an LLM can be asked to
propose a compact directed graph from feature descriptions, domain semantics,
and a statistical fingerprint of the dataset. This should not be confused with
formal causal discovery. In this thesis, the LLM is used as a graph proposal
component: it emits candidate directed edges, the system parses those edges
into an adjacency matrix, and the result is evaluated against an explicit
ground-truth DAG. The scientific question is whether prompt design makes this
proposal process more reliable.

Prompt design is a meaningful research variable because graph recovery
requires more than naming related variables. A model must distinguish direct
from mediated relations, avoid bidirectional or transitive edges unless they
are justified, respect the allowed feature vocabulary, and output a
machine-parseable edge list. Unstructured prompts may produce plausible
correlations, reversed directions, prose explanations, or extra variables.
Structured prompts and prompt decomposition can instead make the model attend
to physical directness, subsystem grouping, and final output constraints.

In this thesis, \emph{prompt decomposition} refers to strategies that divide
the graph-recovery task into smaller internal stages before the final edge
list is produced. The decomposition-oriented strategies are
\texttt{Strategy\_25} (Least-to-Most) and \texttt{Strategy\_26}
(Skeleton-of-Thought). They are compared with \texttt{Strategy\_1}
(Zero-Shot Prompting), \texttt{Strategy\_23} (No-Instruction Control), and
\texttt{Strategy\_24} (Policy-First Structured).

The empirical setting is deliberately focused. The primary evidence is a Tesla
EV result set for two local Ollama models, \texttt{llama3.1:8b} and
\texttt{deepseek-r1:8b}, evaluated on a six-feature energy-consumption
graph-recovery task. This scope makes the work relevant to locally deployed
and infrastructure-constrained fleet analytics, while keeping the claims
auditable. The thesis does not treat the models as inferring true causality
from telemetry alone; it evaluates whether their prompted graph proposals agree
with a predefined physics-informed reference graph.

\section{Research Objective}

The main research question is:

\begin{quote}
To what extent do prompt decomposition and structured prompting strategies
improve LLM-assisted graph recovery from vehicle telematics data when compared
with simpler prompting baselines?
\end{quote}

The question is evaluated in a supervised graph-recovery setting. A predicted
edge set is parsed from each LLM output, converted into an adjacency matrix,
and compared with a fixed Tesla EV ground-truth matrix. The work does not
evaluate formal causal identification, repeated-run self-consistency, or
statistical significance across repeated generations.

The thesis addresses five research questions:

\begin{enumerate}
  \item \textbf{RQ1:} How accurately can \texttt{llama3.1:8b} and
  \texttt{deepseek-r1:8b} recover directed relationships between Tesla EV
  telematics features from a structured dataset summary?

  \item \textbf{RQ2:} How do prompt decomposition strategies, especially
  least-to-most prompting and skeleton-of-thought prompting, compare with
  zero-shot and no-instruction baselines on graph-recovery metrics?

  \item \textbf{RQ3:} How does explicit instruction structuring through a
  policy-first prompt affect precision, recall, false discovery rate, and
  structural graph error?

  \item \textbf{RQ4:} Which evaluated model-strategy combination achieves the
  strongest adjacency-matrix recovery against the Tesla EV ground truth?

  \item \textbf{RQ5:} Which correct, spurious, or reversed relationships
  appear repeatedly across model-strategy outputs, and what do they reveal
  about LLM failure modes in telematics graph recovery?
\end{enumerate}

\section{Contribution Statement}

The thesis makes four contributions:

\begin{itemize}
  \item it defines a reproducible Tesla EV graph-recovery task in which LLM
  outputs are converted into directed adjacency matrices and evaluated against
  a fixed ground-truth DAG;

  \item it compares five prompting strategies across two local open-source
  LLMs, with particular attention to prompt decomposition through
  least-to-most reasoning and skeleton-of-thought prompting;

  \item it reports graph-recovery performance using an eight-metric evaluation
  suite: precision, recall, F1 score, structural Hamming distance, false
  discovery rate, Pearson correlation, Hamming distance, and Frobenius norm;

  \item it analyzes edge-level errors and consensus patterns, showing where
  LLM-assisted graph proposals recover meaningful telemetry relations and
  where they produce false positives or direction reversals.
\end{itemize}

The main empirical finding is that the best observed result is produced by
\texttt{deepseek-r1:8b} with
\strategy{Strategy_26_Skeleton_of_Thought}. In the evaluation matrix, this
combination achieves precision \(0.4444\), recall \(0.6667\), F1 score
\(0.5333\), structural Hamming distance \(7\), false discovery rate
\(0.5556\), Pearson correlation \(0.4303\), Hamming distance \(0.1944\), and
Frobenius norm \(2.6458\). These values support a controlled comparison of
prompt strategies; they do not support claims of general causal validity,
statistical significance, or self-consistency improvement.

\section{Overview of Methodology}

The methodology follows a controlled prompt-to-graph evaluation design. The
Tesla EV dataset is represented through six normalized features:
\texttt{electricenergyremaining}, \texttt{electricremaining},
\texttt{speed\_kmh}, \texttt{acceleration\_kmh2}, \texttt{mileage}, and
\texttt{distance\_m}. A ground-truth directed graph defines six expected
relationships between these features. Each prompt strategy receives the same
feature set and dataset fingerprint, asks the model to output directed edges
in the form \texttt{feature\_a -> feature\_b}, and restricts the final answer
to machine-parseable relationships.

The generated edge list is parsed into a predicted adjacency matrix and
compared with the ground-truth adjacency matrix using the eight reported
metrics. This design keeps the evaluation auditable: every numerical result
in the results chapter is tied to the saved CSV, JSON, and matrix artifacts
for the Tesla EV experiment.

\section{Thesis Structure}

Chapter~\ref{sec:background} reviews the literature on LLM prompting,
structured prompting, prompt decomposition, vehicle telematics analytics, and
graph-based evaluation. Chapter~\ref{sec:task-definition} defines the Tesla
EV dataset, graph-recovery task, ground-truth DAG, and scope of the study.
Chapter~\ref{sec:methodology} presents the methodology, including the
prompt-to-graph pipeline, parser, metrics, and experimental setup.
Chapter~\ref{sec:results} reports the evaluation matrix, model-strategy
comparisons, consensus graph, and edge-level error patterns.
Chapter~\ref{sec:discussion} interprets the results and discusses threats to
validity and future work. Chapter~\ref{sec:conclusion} concludes the thesis.
